\section{Future Work: Virtual Machines}
\label{sec:future}

Containers share the host's kernel. This is what makes \cvexthree records
light enough to run by the hundred on one machine, but it also bounds the
vulnerabilities the format can express. A record cannot install a different
kernel, load kernel modules, change kernel configuration, or run a non-Linux
operating system, and a privileged operation inside a container affects the
host as well. We plan to extend the format with a \emph{VM target}: a
manifest option that replaces the target container with a virtual machine
image, while keeping everything else in the record unchanged (the roles, the
success criteria, the two-phase run protocol, and the dataset layout). The
VM-based generation of \cvex (Section~\ref{sec:background}) already
demonstrated how to provision and instrument such targets; the extension
would bring that capability into the uniform \cvexthree record.

Several classes of vulnerabilities need this extension:

\begin{itemize}
\item \textbf{Kernel vulnerabilities.} Linux kernel flaws such as
CVE-2016-5195 (``Dirty COW''), CVE-2022-0847 (``Dirty Pipe''), and
CVE-2024-1086 (a use-after-free in \texttt{nf\_tables}) depend on a specific
vulnerable kernel version and configuration, which a container cannot
provide.

\item \textbf{Vulnerabilities in non-Linux systems.} Windows vulnerabilities
such as CVE-2017-0144 (the SMBv1 flaw known as EternalBlue), CVE-2019-0708
(``BlueKeep'', in the Remote Desktop service), and CVE-2021-34527
(``PrintNightmare'', in the Print Spooler) require a Windows guest.

\item \textbf{Container and isolation boundaries.} Vulnerabilities in the
container runtime itself, such as CVE-2019-5736 in \texttt{runc}, cannot be
reproduced safely inside the runtime they affect.

\item \textbf{Environment-dependent behavior.} Some user-space
vulnerabilities depend on how the C library and the kernel interact. While
building the corpus we were unable to reproduce CVE-2019-18276, a
\texttt{bash} privilege-management flaw, because it requires a C library and
kernel combination that preserves the saved user ID when \texttt{bash} drops
privileges, which our container environment does not provide.

\item \textbf{Records that need host privileges.} A few records in the
corpus already require additional container capabilities (for example,
permission to mount file systems or to configure traffic control) to
reproduce their vulnerability. Granting these weakens the isolation between
the record and the host; a VM target would restore it.
\end{itemize}

A VM target also enables richer observation. Besides packets and system
calls, a VM-based recorder can capture kernel events, Windows API calls, and
whole-system snapshots taken before and after the attack. The success
criteria of \cvexthree carry over unchanged: a flag readable only by root in
the guest, or a crash witness for a kernel panic, decides success just as it
does for containers.
